\documentclass[11pt]{article}

\usepackage{amsmath,amssymb,graphicx}
\usepackage[margin=1in]{geometry}
\usepackage{booktabs}

\title{Dirac--K\"ahler Braid Topology and Protection Structure on a Kernel-Induced Cochain Operator Architecture}
\author{Tomislav Rupi\'c \\ Independent Researcher}
\date{March 2026}

\begin{document}

\maketitle

\begin{abstract}
The first Phase III release established that the discrete Dirac--K\"ahler sector changes collision phenomenology through structured grade exchange, but does not yet produce a persistence-level advantage over the matched scalar control. The next question is therefore narrower and more structural: does the surviving clustered Dirac--K\"ahler family support a reproducible local topology, and if so, what protects or destroys it? Stages 23.3 and 23.4 first close the most direct persistence extensions. Weak internal grade coupling, phase-lock constraints, curvature proxies, and small geometry variations do not convert the clustered encounter into a bounded composite or a metastable transport family. Stages 23.5 through 23.8 then shift from persistence to topology. A 9-run energy-flow map finds a braid-like exchange class in $8/9$ local perturbations. A 9-run phase scan shows that the braid family survives across a finite corridor, with only a localized smeared midpoint window. A 9-run notch scan resolves that midpoint instability as a broad band rather than a single tuned point. Finally, a 12-run mechanism scan identifies a clear protection structure: phase is edge-sensitive, width symmetry is symmetry-sensitive, weak anisotropy remains topology-preserving, and grade balance is fragile under a zero-form-heavy perturbation. The bounded conclusion is that clustered Dirac--K\"ahler encounters support a real braid-like topology family with an explicit local protection map, while no claim is made for binding, confinement, or emergent particles.
\end{abstract}

\section{Why a Second Phase III Paper Exists}

Paper 23.1 opened the Dirac--K\"ahler sector conservatively. The result was positive only in the narrow sense that graded first-order propagation creates a new collision-texture layer. It was not positive on the stronger axis first hoped for: no bounded composite regime appeared, the matched scalar control remained competitive or stronger on persistence, and the best clustered anomaly did not survive refinement as a binding candidate.

That leaves one lawful continuation. If the clustered Dirac--K\"ahler family is not yet a persistence family, then the next question is whether it is at least a structured topology family. This changes the target from ``does it bind?'' to ``does it preserve a reproducible local exchange class under controlled perturbation?'' Stages 23.3 through 23.8 answer exactly that question and nothing broader.

\section{Negative Boundary Before Topology}

\subsection{Stage 23.3: Stabilization Triggers Fail}

Stage 23.3 takes the single refinement-stable clustered mixed encounter from Paper 23.1 and applies one weak stabilizing trigger at a time. The probes are internal grade coupling, phase-lock composite constraints, and a self-induced curvature proxy. The scan remains fully negative on the persistence axis:
\begin{itemize}
\item $9/9$ runs remain unresolved / mixed,
\item $9/9$ runs remain only weak persistence gain,
\item $0$ promotion candidates appear,
\item $0$ bounded oscillatory separations appear.
\end{itemize}

The important point is not inactivity but insufficiency. Composite lifetime remains high at $0.7517$, binding persistence remains $1.0$, and grade transfer remains nonzero across the branch, but none of the weak local extensions converts the event into a bounded composite regime. Phase locking can improve coherence, yet coherence alone is not capture.

\subsection{Stage 23.4: Metastable Transport Does Not Open}

Stage 23.4 tests the next nearby possibility: perhaps the clustered family does not bind, but does open into a metastable transport window under small changes of separation, width contrast, or phase relation. It does not. Across 9 runs:
\begin{itemize}
\item all runs remain localized encounters,
\item no stable metastable-transport triplets appear,
\item no proto-bound composite appears.
\end{itemize}

These two negative stages matter because they remove the simplest persistence interpretations. The clustered family is real, but its reality is not persistence-level. That clears the way for a topology-first read.

\section{Stage 23.5: A Local Braid-Like Exchange Family Appears}

Stage 23.5 asks whether the clustered Dirac--K\"ahler encounter carries a reproducible local energy-flow topology under small phase, width, and separation perturbations. The answer is the first clear exploratory positive in this subprogram. The 9-run map returns:
\begin{itemize}
\item $8/9$ braid-like exchange,
\item $1/9$ unresolved mixed topology,
\item stable topology families under width variation and separation variation,
\item transfer-dominant secondary labeling in all $9$ runs.
\end{itemize}

This matters for two reasons. First, the clustered family does not dissolve into diffuse washout once local perturbations are introduced. Second, the same braid-like exchange class reappears across more than one perturbation axis, which means the phenomenon is not tied to a single tuned geometry. The claim remains tightly bounded: the result is about local topology class, not about binding or protected matter.

\section{Stages 23.6 and 23.7: Phase Corridor and Midpoint Band}

\subsection{Stage 23.6: Finite Corridor, Not Monotonic Collapse}

Once Stage 23.5 shows a candidate braid-like family, phase becomes the next natural axis to test. Stage 23.6 runs a 9-case phase-sensitivity scan with mild local skew and weak grade-coupling modulation checks. The result is again structured rather than monotonic:
\begin{itemize}
\item $8/9$ runs are braid-like exchange,
\item $1/9$ run is transfer-smeared,
\item baseline, quarter-offset, three-quarter-offset, and anti-phase anchors remain braid-like,
\item skew and weak $\beta$ perturbations also remain braid-like.
\end{itemize}

The only failure occurs at the midpoint detuning window. This is important because it shows that phase sensitivity exists, but does not act like a simple global degradation parameter. The braid family occupies a finite corridor.

\subsection{Stage 23.7: The Midpoint Instability Resolves as a Broad Band}

Stage 23.7 resolves the suspicious midpoint window with a local 9-run notch scan centered on phase-offset fraction $0.5$. The resulting sequence is explicit:
\[
0.375 \rightarrow \text{braid}, \quad
0.425 \rightarrow \text{braid}, \quad
0.45 \rightarrow \text{braid}, \quad
0.475 \rightarrow \text{braid},
\]
\[
0.5 \rightarrow \text{smeared}, \quad
0.525 \rightarrow \text{smeared}, \quad
0.55 \rightarrow \text{smeared}, \quad
0.575 \rightarrow \text{smeared}, \quad
0.625 \rightarrow \text{braid}.
\]

In count form the scan gives $5/9$ braid-like exchange and $4/9$ transfer-smeared, with the midpoint corridor interpreted as broad rather than sharp. The upper-anchor recovery at $0.625$ is decisive. It shows that the system does not simply degrade as phase detuning increases. Instead, the clustered Dirac--K\"ahler family passes through a finite smeared band and then recovers the braid class outside it.

\section{Stage 23.8: Mechanism Scan and Protection Structure}

Stage 23.8 asks the next honest question. If the braid family is real, what preserves it, and what breaks it? The scan uses four isolated branches with three runs each:
\begin{itemize}
\item grade-balance perturbation,
\item phase-band edge protection,
\item width / symmetry breaking,
\item weak anisotropy breaking.
\end{itemize}

The 12-run result is positive in the sense required by the prompt: at least one clear protection or breaking rule is identified. In fact, the scan yields three clean rules and one fragile axis. The topological counts are:
\begin{itemize}
\item $8/12$ braid-like exchange,
\item $4/12$ transfer-smeared.
\end{itemize}

The key structural feature is that the smeared cases are not diffuse collapse events. Even the transfer-smeared runs retain high flow concentration, between $0.8667$ and $0.8928$, and grade-exchange coherence remains tightly bounded between $0.4688$ and $0.4746$. The topology transition therefore occurs within an organized collision family, not through simple washout.

\subsection{Grade Balance Is Fragile}

The grade-balance branch holds geometry and phase fixed while changing the initial $0$-form to $1$-form support ratio. The balanced baseline is braid-like. The one-form-enhanced run also remains braid-like. But the zero-form-enhanced run becomes transfer-smeared. The protection read is therefore not ``grade balance fully protects'' but rather ``grade imbalance can destabilize the family,'' with the tested zero-form-heavy perturbation acting as a failure mode.

\subsection{Phase Is Edge-Sensitive}

The phase-edge branch is the cleanest branch in the paper. The lower braid anchor at phase-offset fraction $0.475$ is braid-like, the midpoint sample at $0.5$ is transfer-smeared, and the upper anchor at $0.625$ returns to braid-like exchange. This is a reproducible braid-to-smeared-to-braid pattern, so the correct label is edge-sensitive. Phase is not a generic damage knob. It defines a local band edge.

\subsection{Width Symmetry Is Sensitivity, Not Decoration}

The symmetry branch fixes phase at the lower braid anchor and introduces width asymmetry. The symmetric baseline is braid-like. Both asymmetric cases, Packet A wider and Packet B wider, become transfer-smeared. This is strong evidence that local geometric symmetry is part of the topology support structure. The result is not just that geometry matters. It is that the braid label is specifically symmetry-sensitive under the tested width imbalance.

\subsection{Weak Anisotropy Preserves the Braid Class}

The anisotropy branch gives the most stable result. The isotropic baseline, weak $x$-stretch, and weak $y$-stretch all remain braid-like exchange. This is the clean positive protection read in the scan: mild directional anisotropy does not destroy the topology class at the tested strengths.

\begin{table}[h]
\centering
\scriptsize
\begin{tabular}{@{}p{1.2in}p{1.8in}p{1.2in}p{1.3in}@{}}
\toprule
Branch & Three-run outcome & Protection label & Read \\
\midrule
Grade balance & braid, smeared, braid & fragile & zero-form-heavy bias destabilizes \\
Phase edge & braid, smeared, braid & edge-sensitive & finite braid-to-smeared band edge \\
Width symmetry & braid, smeared, smeared & symmetry-sensitive & width asymmetry breaks the braid class \\
Weak anisotropy & braid, braid, braid & protected & weak directional bias is tolerated \\
\bottomrule
\end{tabular}
\caption{Stage 23.8 mechanism scan. The topology family has a real local protection map rather than a single tuned point.}
\end{table}

\section{What Paper 23.2 Adds}

Paper 23.1 showed that the Dirac--K\"ahler sector changes collision phenomenology. Paper 23.2 sharpens that statement in a more specific direction.

\subsection{The Positive Signal Is Topological, Not Persistent}

Stages 23.3 and 23.4 are necessary because they block the easy misread. The clustered family does not become a bounded composite under weak stabilizing extensions and does not open into metastable transport under nearby geometric perturbations. The positive result therefore has to live elsewhere.

\subsection{The Clustered Family Supports a Reproducible Exchange Topology}

Stages 23.5 through 23.7 show that the surviving clustered family is not a one-off mixed encounter. It supports a braid-like local exchange class, preserves that class across width and separation perturbations, survives across a finite phase corridor, and resolves its midpoint failure as a broad internal phase band instead of a single accidental point.

\subsection{The Topology Family Has a Real Protection Map}

Stage 23.8 turns the exploratory positive into a structural positive. The braid family is not equally sensitive to all perturbations. Phase behaves like an edge parameter, width symmetry behaves like a fragile support condition, weak anisotropy is tolerated, and grade imbalance can destabilize the class. That is exactly the kind of result that justifies the term ``protection structure'' in the bounded operational sense used here.

\begin{table}[h]
\centering
\scriptsize
\begin{tabular}{@{}p{0.7in}p{0.5in}p{1.2in}p{2.7in}@{}}
\toprule
Stage & Runs & Result type & Bounded conclusion \\
\midrule
23.3 & 9 & negative pilot & weak local triggers do not create bounded composites \\
23.4 & 9 & negative boundary & nearby geometric variation does not open metastable transport \\
23.5 & 9 & exploratory positive & clustered DK encounters support a braid-like local topology class \\
23.6 & 9 & corridor positive & the braid family survives across a finite phase corridor \\
23.7 & 9 & band positive & the midpoint instability resolves as a broad smeared band \\
23.8 & 12 & mechanism positive & the braid family exhibits an explicit local protection map \\
\bottomrule
\end{tabular}
\caption{The Paper 23.2 arc. The persistence line closes before the topology line opens.}
\end{table}

\section{Conclusion}

The correct claim for this second Phase III paper is narrower than binding but stronger than the opener claim. The clustered Dirac--K\"ahler family does not merely exhibit nontrivial collision texture. It supports a reproducible braid-like exchange topology with a mapped local protection structure.

That statement remains bounded in four ways.
\begin{itemize}
\item No bound state is claimed.
\item No confinement or particle interpretation is claimed.
\item The protection language refers only to recoverable topology class under local perturbation.
\item The result is local to the tested clustered Dirac--K\"ahler family on the frozen periodic $2$D complex.
\end{itemize}

Within those bounds the result is real. Phase III now contains two distinct positive statements:
\begin{itemize}
\item Paper 23.1: the Dirac--K\"ahler sector opens a new collision-texture layer through structured grade exchange.
\item Paper 23.2: the clustered Dirac--K\"ahler family supports a braid-like topology class with an explicit local protection map.
\end{itemize}

That is enough to justify the next continuation of the program. The lawful move is no longer a blind search for binding. It is a focused study of why this topology family exists, how far its protection extends under refinement and operator variation, and whether the protection map survives beyond the minimal clustered construction.

\end{document}
